\section{The single-file formalization and provenance}\label{sec:formalization}

The companion file \texttt{Erdos12AllParts.lean} imports only Mathlib.
It contains the local dependency closure of the reciprocal-sum argument
and the two attributed AlphaProof Nexus proof developments. Its public
interface is as follows.

\begin{center}
\begin{tabular}{p{0.22\textwidth}p{0.69\textwidth}}
\hline
Entry point & Mathematical statement\\
\hline
\texttt{part\_i} & A good infinite set has positive lower limit at square-root scale.\\[3pt]
\texttt{part\_ii} & The universal power-saving assertion is false.\\[3pt]
\texttt{part\_iii} & Every good infinite set has a convergent reciprocal sum.\\[3pt]
\texttt{all\_parts} & The conjunction of the three preceding answers.\\[3pt]
\texttt{parameterwise\_}\newline\texttt{dense} & The set and eventual counting threshold may depend on the exponent.\\[3pt]
\texttt{uniform\_}\newline\texttt{reciprocal\_bound} & One constant bounds the reciprocal sum of every set with Property P.\\
\hline
\end{tabular}
\end{center}
All entry points belong to the namespace \texttt{Erdos12AllParts}.
The definition \texttt{IsGood} retains the quantifiers and strict
inequalities of the original Formal Conjectures statement.

The check uses Lean \texttt{4.35.0-rc2} and Mathlib commit
\nolinkurl{065356127b1dc0016f66b7283ce0ce2c4055aa55}.
The prior proof files were taken from Google DeepMind's
\texttt{alphaproof-nexus-results} repository at commit
\nolinkurl{0647711a71183c1ea492ad60860776617ce1ea88}.
They are distributed under the Apache License, Version 2.0, with
copyright 2025 Google LLC. The companion file retains these notices and
the license. Their original environment used Lean 4.27.0;
the adaptation isolates their namespaces, removes redundant options and
imports, updates the image-cardinality API, and expands the
interval-residue counting lemma into a direct proof.

The analytic formalization also retains the Euler--Maclaurin and
Mertens developments from \texttt{PrimeNumberTheoremAnd}, obtained via
\texttt{chen\_theorem}, and the vendored large-sieve development from
\texttt{goldbach-lean} \cite{PNTCode,GoldbachCode}. Their recorded source
revisions and attributions are preserved in the Lean header and the
validation/provenance files; these components are likewise recorded under
Apache License 2.0 in the supplied project.

The final checks compile the combined source, check each public entry
point, and inspect its axiom dependencies. The only axioms reported are
\texttt{propext}, \texttt{Classical.choice}, and \texttt{Quot.sound}.
The file introduces no admitted theorem and does not obtain the first
two results by importing placeholder catalog declarations. A separate
check compares all three endpoints with the pinned original catalog
statements. The validation receipt records the source digest and commands.

\section*{Acknowledgments and contribution statement}

Parts (i) and (ii), including the constructions and the released formal
proofs adapted here, are credited to the AlphaProof Nexus work of
Google DeepMind and Tsoukalas, Kovsharov, Shirobokov et al.\ \cite{TKS}.
The present manuscript supplies a unified exposition and the recorded
adaptation to the pinned Lean environment. The four named authors retain
responsibility for the reciprocal-sum argument, the integration, and the
final mathematical claims.

ChatGPT and Codex, developed by OpenAI, assisted with the mathematical
development of the reciprocal-sum argument, checking calculations,
literature searches, source integration, and manuscript preparation.
The earlier reductions and analytic ingredients used in part (iii)
are credited in the introduction and body of the proof.

\begin{thebibliography}{99}

\bibitem{AB}
B. Alexeev, K. Barreto, Y. Li, J. D. Lichtman, L. Price,
J. I. Shah, Q. Tang and T. Tao,
\emph{Primitive sets and von Mangoldt chains: Erd\H{o}s Problem
\#1196 and beyond},
preprint (2026),
\href{https://arxiv.org/abs/2605.00301v1}{arXiv:2605.00301v1}.

\bibitem{BV}
T. Bera and G. K. Viswanadham,
\emph{Restriction estimates with sifted integers},
preprint (2026),
\href{https://arxiv.org/abs/2512.21640v3}{arXiv:2512.21640v3}.

\bibitem{ES}
P. Erd\H{o}s and A. S\'ark\"ozy,
\emph{On the divisibility properties of sequences of integers},
Proc. London Math. Soc. (3) \textbf{21} (1970), 97--101.
\href{https://doi.org/10.1112/plms/s3-21.1.97}{doi:10.1112/plms/s3-21.1.97}.

\bibitem{NexusI}
Google DeepMind,
\emph{AlphaProof Nexus results: Erd\H{o}s Problem 12(i)},
Lean source (2026), Apache License, Version 2.0.
Repository commit \nolinkurl{0647711a71183c1ea492ad60860776617ce1ea88}.
\url{https://github.com/google-deepmind/alphaproof-nexus-results/blob/0647711a71183c1ea492ad60860776617ce1ea88/APNOutputs/ErdosProblems/erdos_12.parts.i.lean}.

\bibitem{NexusII}
Google DeepMind,
\emph{AlphaProof Nexus results: Erd\H{o}s Problem 12(ii)},
Lean source (2026), Apache License, Version 2.0.
Repository commit \nolinkurl{0647711a71183c1ea492ad60860776617ce1ea88}.
\url{https://github.com/google-deepmind/alphaproof-nexus-results/blob/0647711a71183c1ea492ad60860776617ce1ea88/APNOutputs/ErdosProblems/erdos_12.parts.ii.lean}.

\bibitem{MV}
H. L. Montgomery and R. C. Vaughan,
\emph{\href{https://doi.org/10.1112/S0025579300004708}{The large sieve}},
Mathematika \textbf{20} (1973), 119--134.

\bibitem{MVbook}
H. L. Montgomery and R. C. Vaughan,
\emph{Multiplicative Number Theory I. Classical Theory},
Cambridge Studies in Advanced Mathematics, vol.~97,
Cambridge University Press, Cambridge, 2007.

\bibitem{Plasma}
Plasma AI,
\emph{Composite-modulus large sieve for Property P},
research notes, Erd\H{o}s Problems (2026),
\href{https://erdosproblems.ai/wiki/research/leads/property_p_sieve_multiplicity}{online edition},
accessed 4 November 2026.

\bibitem{PNTCode}
\emph{PrimeNumberTheoremAnd}, formalization project,
\url{https://github.com/AlexKontorovich/PrimeNumberTheoremAnd},
commit \nolinkurl{c6c73610b406689c4b58325cbe1342ecca25d755};
retrieved through \url{https://github.com/lixiang90/chen_theorem},
commit \nolinkurl{11dcca8fc3ab340358503278c46d72aef072a6e7}.

\bibitem{GoldbachCode}
\emph{goldbach-lean}, large-sieve formalization,
\url{https://github.com/subfish-zhou/goldbach-lean},
commit \nolinkurl{9d4c3a4010d42aedd9e7231bf94386019cfffccc}.

\bibitem{TKS}
G. Tsoukalas, A. Kovsharov, S. Shirobokov et al.,
\emph{Advancing mathematics research with AI-driven formal
proof search},
preprint (2026),
\href{https://arxiv.org/abs/2605.22763v2}{arXiv:2605.22763v2},
Appendix~B.4.

\end{thebibliography}
